\documentclass[11pt]{article}

% --- Preamble: restricted to packages confirmed present in this TeX Live 2025
%     basic install (siunitx and enumitem are NOT available here). ------------
\usepackage[margin=1in]{geometry}
\usepackage{amsmath,amssymb}
\usepackage{graphicx}
\usepackage{booktabs}
\usepackage{xcolor}
\usepackage{caption}
\usepackage{listings}
\usepackage[hidelinks]{hyperref}

% Robust verbatim-in-text for names with underscores (avoids manual escaping).
\newcommand{\code}[1]{\texttt{\detokenize{#1}}}
\lstset{basicstyle=\ttfamily\footnotesize,breaklines=true,frame=single,
        columns=fullflexible,xleftmargin=1.5em}

\title{\textbf{Continuum Mechanics of Catholyte Contact in Solid-State\\
Sodium-Ion Batteries: A Predictive Design Framework}\\[4pt]
\large Modeling support for the plastic-glassy NACS catholyte manuscript}
\author{Continuum modeling, ecmLab}
\date{9 June 2026 \quad(\emph{working report --- status summary})}

\begin{document}
\maketitle

\begin{abstract}
\noindent
This report summarizes the continuum finite-element modeling built to support
the manuscript \emph{``Plastic Glassy Catholyte Unlocking Stable Solid-State
Sodium-Ion Batteries''} (Gan \textit{et al.}). The manuscript is wholly
experimental; the modeling described here is a \emph{new, theory-first
predictive framework} intended to open the paper. Its purpose is not to
reproduce a measurement but to answer a design question: \emph{what mechanical
properties must a catholyte possess to maintain particle-to-particle contact in
a composite cathode while minimizing the required stack pressure?} The model
treats the cathode-active-material (CAM) / solid-electrolyte (SE) interface with
two complementary descriptions --- node-to-segment contact and a cohesive-zone
model (CZM) --- and sweeps the catholyte's mechanical properties to produce a
\emph{design map}. The conclusions then feed a downstream DFT screening step
that places candidate materials on that map. The current production model uses
rate-independent $J_2$ plasticity for the catholyte and isotropic linear
elasticity for the NVP cathode; a temperature-dependent elasto-viscoplastic
catholyte law was also developed and is parked pending experimental confirmation
of the flow behavior. We report the modeling rationale, both constitutive
options, the basis borrowed from a closely related published study, and the
implementation status (Phases~0--1 complete).
\end{abstract}

\section{Objective and scientific context}

The central mechanical thesis of the manuscript is that a \emph{rigid} catholyte
cannot follow the CAM's volume change during cycling, so it delaminates and
loses contact, forcing the use of high stack pressure ($>150$\,MPa). The
proposed NACS (Na$_2$Al$_{1.35}$Cl$_{4.05}$S) glass instead undergoes a low-endothermic
\emph{plastic phase transition} on in-situ heating and flows as a viscous,
binder-like catholyte, maintaining intimate contact at $\sim$10\,MPa or even
pressure-free assembly. The experimental evidence (nanoindentation, DSC, SEM,
long-term cycling) is strong but qualitative on the mechanics.

The modeling goal is therefore to provide the \textbf{predictive trend}:
quantify interface contact loss and the minimum stack pressure required to
preserve contact, as a function of the catholyte's stiffness, hardness and
viscous (rate-dependent) character. Because the intent is a transferable design
map rather than a microstructure reconstruction, an idealized single-particle
representative-interface geometry and a property sweep are the appropriate level
of fidelity.

\section{Physical problem and modeling strategy}

\paragraph{Geometry and unit system.} A 2-D $5\times5$\,$\mu$m representative
cell (unit system $\mu$m--MPa--$\mu$N; $1$\,MPa${=}1\,\mu$N$/\mu$m$^2$): an NVP
cathode quarter disc (radius $R\approx3.99$\,$\mu$m, $\approx$50\% area fraction)
embedded in an L-shaped NACS catholyte domain --- i.e.\ a $\mu$m-scale NVP
active particle. The shared quarter-circle arc is the interface where all reported
quantities live. The bulk mechanics is scale-invariant under uniform geometric
scaling, so the $\mu$m setting matters chiefly by fixing the physical
cohesive-zone length that governs CZM delamination. The mesh is graded ---
element edge $0.02$\,$\mu$m on the arc, coarsening to $0.20$\,$\mu$m in the bulk
--- and fully
recombined to QUAD4 ($\sim$10{,}050 elements, 313 elements along the arc).

\paragraph{Two complementary interface models.} Both share the geometry,
loading, and bulk materials; they differ only in how the arc is treated
(Table~\ref{tab:models}).

\begin{table}[h]
\centering
\caption{The two complementary interface descriptions.}
\label{tab:models}
\begin{tabular}{lll}
\toprule
 & \code{1_contact.i} & \code{2_czm.i} \\
\midrule
Interface & non-conformal (two arcs) & conformal + \code{BreakMeshByBlockGenerator} \\
Physics   & penalty node-to-segment contact & bilinear mixed-mode CZM \\
Mesh      & \code{mesh_5050.msh} & \code{mesh_czm_5050.msh} \\
Reports   & gap / contact loss & traction, jump, damage (delamination) \\
\bottomrule
\end{tabular}
\end{table}

\paragraph{Loading.} The CAM volume change is imposed as a thermal-expansion
eigenstrain proxy (a charge/discharge ramp, $\sim$8\% volumetric swelling at
peak), superposed on a stack-pressure boundary condition ramped on the top face.
The NVP cathode is treated as isotropic linear elastic (no plasticity); the
catholyte carries the inelastic (plastic) response.

\section{Catholyte constitutive model}

\subsection{Catholyte: $J_2$ plasticity (current) and a parked viscoplastic extension}

The \textbf{current production model} is rate-independent $J_2$ plasticity:
linear elasticity plus an isotropic (von Mises) yield whose strength is tied to
the measured Vickers hardness through the Tabor relation,
\begin{equation}
  \sigma_y \;\approx\; \frac{H_v}{3}.
  \label{eq:tabor}
\end{equation}
This captures the onset of plastic flow but \emph{not} the time/rate-dependent
viscous flow that is the essence of the plasticized NACS ``binder.'' A
\textbf{rate-dependent viscoplastic extension} was therefore also developed ---
linear elasticity $+$ static $J_2$ yield (Eq.~\ref{eq:tabor}) $+$ Norton
(power-law) creep,
\begin{equation}
  \dot{\varepsilon}^{\,\mathrm{cr}} \;=\; A\,\sigma_{\mathrm{eq}}^{\,n},
  \label{eq:norton}
\end{equation}
implemented as \code{PowerLawCreepStressUpdate} alongside
\code{IsotropicPlasticityStressUpdate} in \code{ComputeMultipleInelasticStress}.
Equation~\ref{eq:norton} is the MOOSE-native analog of the rate-dependent yield
surface used in the reference study (Sec.~\ref{sec:ref}),
\begin{equation}
  \frac{\bar{\sigma}}{Y} \;=\;
  \left(1+\frac{\dot{\varepsilon}^{\,p}}{\dot{\varepsilon}_0}\right)^{1/n},
  \label{eq:cowper}
\end{equation}
recovering stiff behavior at low temperature and viscous flow as the glass
softens. \textbf{This viscoplastic version is parked in \code{viscoplastic/}}
pending experimental confirmation of the flow behavior; the production inputs
currently use $J_2$ only.

\subsection{Assumed parameters (placeholders)}

Pending the redone hold-segment nanoindentation, we adopt the
\textbf{placeholder} temperature presets in Table~\ref{tab:presets}. $E$ and
$H_v$ are the $J_2$ (and DFT-screening) properties; the creep coefficient $A$
(with $n=4$, mirroring the Li$_6$PS$_5$Cl reference) belongs to the \emph{parked}
viscoplastic version, where it increases by roughly two orders per step so the
response evolves from rigid (25\,$^\circ$C) through softening (60\,$^\circ$C) to
fluid-like (90\,$^\circ$C). Values are overridden on the command line to sweep
and build the design map.

\begin{table}[h]
\centering
\caption{Assumed (placeholder) catholyte properties --- \emph{to be recalibrated
to the redone nanoindentation}. $\nu=0.30$ throughout.}
\label{tab:presets}
\begin{tabular}{lcccc}
\toprule
State & $T$ ($^\circ$C) & $E$ (MPa) & $H_v$ (MPa) & $A$ (Norton, $n{=}4$) \\
\midrule
rigid glass     & 25 & 1000 & 30 & $1\times10^{-7}$ \\
softening       & 60 &  500 & 12 & $1\times10^{-5}$ \\
plastic/slurry  & 90 &  100 &  3 & $1\times10^{-3}$ \\
\bottomrule
\end{tabular}
\end{table}

\noindent\textbf{Important caveat.} The absolute value of $A$ is tied to the
(currently proxy) loading timescale; it must be rescaled once the real
charge/discharge rate is set so that creep accumulates an $O(1\%)$ strain over a
cycle. The temperature dependence is presently baked into the preset (no
coupling to the swelling-proxy temperature field), which cleanly avoids
overloading that field with two physical meanings.

\section{Basis in the literature and our extension}
\label{sec:ref}

The modeling approach closely follows Shi \textit{et al.}, \emph{``Stack
pressure effects and viscoplastic deformation in argyrodite solid-state
electrolyte,''} \emph{Matter} \textbf{9}, 102767 (2026), which provides a
near-turnkey blueprint. We borrow:

\begin{itemize}
  \item \textbf{Constitutive form} --- the rate-dependent yield surface,
  Eq.~\ref{eq:cowper}. Their Li$_6$PS$_5$Cl parameters
  ($E=22.1$\,GPa, $\nu=0.30$, $Y=0.35$\,GPa, $\dot{\varepsilon}_0=10^{-5}$\,s$^{-1}$,
  $n=4$) anchor the \emph{rigid} end of our sweep.
  \item \textbf{Calibration workflow} --- inverse-FEM fitting of $(E,Y,\dot{\varepsilon}_0,n)$
  to hold-segment nanoindentation, using a \emph{relative}-depth-change scheme
  robust to porous/rough surfaces.
  \item \textbf{Design-map template} --- their contact-ratio versus
  \emph{normalized} stack pressure $p/Y_0$, which validates $p/Y$ as the
  governing dimensionless group and is the direct template for our opening
  figure.
\end{itemize}

\noindent\textbf{Our novelty.} Their model is isothermal. The plastic phase
transition of NACS is precisely a \emph{temperature-activated} change in the
viscoplastic parameters $E(T)$, $Y(T)$, $A(T)$ across the glass transition ---
so the temperature-dependent constitutive law (Table~\ref{tab:presets}) and the
resulting temperature axis on the design map are a genuine extension, not a
re-run.

\section{Design-map property axes and the DFT handoff}

The screening axes are chosen to be both physically governing \emph{and}
computable by DFT, since DFT is the downstream filter that ranks candidate
catholytes (Table~\ref{tab:axes}).

\begin{table}[h]
\centering
\caption{Design-map / screening axes and their DFT route.}
\label{tab:axes}
\begin{tabular}{lll}
\toprule
Property & Role & DFT route \\
\midrule
$E$ (or $E/\sigma_y$) & elastic conformability & elastic tensor \\
$H_v \to \sigma_y$ & plastic gap-filling & $H_v(G,B)$ models, then Eq.~\ref{eq:tabor} \\
adhesion $G_{Ic}$ & resistance to delamination & interface/surface energy \\
Pugh ratio $G/B$ & ductility (optional) & elastic tensor \\
viscosity $A$ & viscous flow rate & \emph{not} DFT-screenable; a model knob \\
\bottomrule
\end{tabular}
\end{table}

\paragraph{Hardness from DFT.} DFT does not yield $H_v$ directly, but
well-established models map the computed elastic moduli to it. With shear modulus
$G$, bulk modulus $B$, and Pugh's ratio $k=G/B$:
\begin{align}
  H_v &= 2\,(k^2 G)^{0.585} - 3 && \text{(Chen--Niu)}, \\
  H_v &= 0.92\,k^{1.137} G^{0.708} && \text{(Tian)}, \\
  H_v &\approx 0.151\,G && \text{(Teter)}.
\end{align}
Thus $H_v \approx f(G,\,G/B)$: hardness scales with the magnitude of $G$, while
Pugh's ratio supplies the ductility correction --- the two are intimately
related but not identical. A single DFT elastic calculation per candidate yields
$\{E,\nu,G,B\}$ and hence \emph{both} primary map axes ($E$ directly, and
$\sigma_y=H_v/3$). These models are calibrated on hard crystalline solids, so
they should be used for \emph{relative} ranking; a porosity knock-down is needed
to compare dense-cell DFT moduli with the cold-pressed porous pellet, and an
amorphous NACS requires a melt-quench supercell with Voigt--Reuss--Hill
averaging. Because for a real material $E$ and $H_v$ both derive from the same
$\{G,B\}$, candidates populate only a physically allowed band of the $(E,H_v)$
plane --- itself useful information about which corner of the map is reachable.

\paragraph{Pipeline.}
\begin{center}
continuum design map $\;\longrightarrow\;$ DFT computes $\{E,H_v,G_{Ic}\}$ per
candidate $\;\longrightarrow\;$ overlay on the map $\;\longrightarrow\;$ select the
material deepest in the low-loss / low-pressure region.
\end{center}
The DFT work proceeds in a parallel repository
(\code{abInit/dftecm/paper_yangzhao_catholyteMechanics}); elastic results are
pending.

\section{Implementation status}

\subsection{Phase 0 --- housekeeping (complete)}
\begin{itemize}
  \item Both meshes regenerated via \code{make_mesh.py} (pure QUAD4; interface
  sidesets verified: \code{block_NVP_right}/\code{block_NACS_left} for contact,
  \code{block_NVP_block_NACS} for the CZM split).
  \item Sweep scripts de-staled: executable \code{contact_loss-opt} $\to$
  \code{ecm-opt} (app name \code{ecm}) via an \code{EXE} variable; HPC
  \code{BASE} path parameterized.
  \item Both inputs smoke-tested locally with \code{ecm-opt}.
\end{itemize}

\subsection{Phase 1 --- catholyte model and elastic cathode (complete)}
Both inputs run with the $J_2$ catholyte and an isotropic-elastic NVP cathode
(\code{ComputeFiniteStrainElasticStress}), and converge. The parked viscoplastic
catholyte (in \code{viscoplastic/}) adds a Norton-creep material block:
\begin{lstlisting}
[./power_law_creep_NACS]
  type = PowerLawCreepStressUpdate
  block = 'block_NACS'
  coefficient = ${creep_coef}   # Norton A (ASSUMED; recalibrate)
  n_exponent  = ${creep_exp}    # Norton n (ASSUMED)
  activation_energy = 0         # T-dependence baked into A (no temp coupling)
  max_inelastic_increment = 1e-3
[../]
[./radial_return_stress_NACS]
  type = ComputeMultipleInelasticStress
  tangent_operator = nonlinear
  inelastic_models = 'isotropic_plasticity_NACS power_law_creep_NACS'
  block = 'block_NACS'
[../]
\end{lstlisting}

\paragraph{Verification (parked viscoplastic version).} It was confirmed active
by three independent checks (CZM case, default mesh):
\begin{itemize}
  \item \emph{Residual signature.} Switching on creep raises the Newton-iteration
  count of a representative step from 3 (rate-independent) to 18 --- the creep
  term genuinely enters the residual.
  \item \emph{Stress relaxation} (the binder mechanism). At $t=0.16$ the peak von
  Mises stress in the catholyte drops from $4.000$\,MPa (creep off) to
  $3.899$\,MPa at $A=10^{-3}$ --- a $\sim$2.5\% relaxation, the textbook signature
  of viscous flow relieving stress.
  \item \emph{Gap-controlled kinematics (so far).} The arc-apex gap is nearly
  unchanged ($4.3215\times10^{-2}$ off vs.\ $4.3220\times10^{-2}$ at $A=10^{-3}$):
  early in a single load the response is gap-controlled and the low yield
  ($\sigma_y\!\approx\!3.3$\,MPa) caps the driving stress, so the
  contact-retention payoff of viscous flow becomes first order only under cyclic
  loading and after the rate is recalibrated (Phases~2--3).
\end{itemize}

\subsection{Convergence and performance}
Profiling identified the cost driver and the fixes (all verified locally):
\begin{itemize}
  \item \textbf{Diagnosis.} The hand-coded CZM / creep / penalty-contact tangents
  are inexact, so Newton converges only linearly --- the residual reaches
  $\sim$$10^{-7}$ in 2--3 iterations then \emph{crawls} to the former
  \code{nl_abs_tol} of $10^{-8}$. Near peak swelling the CZM cohesive-softening
  (indefinite) tangent additionally makes Newton oscillate and diverge.
  \item \textbf{Tolerance (baked in).} Loosening \code{nl_abs_tol}
  $10^{-8}\!\to\!10^{-7}$ drops the wasted crawl: \textbf{2.4$\times$} fewer
  iterations for the CZM input and \textbf{3.8$\times$} for the contact input,
  with gap/traction unchanged to $\sim$$10^{-5}$ relative. A backtracking line
  search does not help and \emph{worsens} the CZM peak-load case.
  \item \textbf{Exact Jacobian via AD.} The full bulk chain (elasticity,
  plasticity, Norton creep, eigenstrain) has AD versions; converting it makes the
  bulk tangent exact and Newton truly quadratic --- \textbf{max 3 iterations per
  step}, and the contact model runs \emph{through peak swelling with zero
  cutbacks} (vs.\ 9 cutbacks / 99-iteration stalls without it). Two caveats:
  \code{ADBiLinearMixedModeTraction} does not exist, so the CZM interface tangent
  stays hand-coded; and each AD assembly is $\sim$$10\times$ costlier, so AD is
  $\sim$$6\times$ slower in \emph{wall-clock} despite fewer iterations. AD is thus
  kept as a robustness tool, not the production default.
  \item \textbf{Adaptive timestep.} \code{IterationAdaptiveDT} was clamped to
  $\Delta t\le0.02$ by 50 dense output sync times; output is now decoupled from
  stepping (sync only at the critical times 0.05, 0.5, 1.0; \code{optimal_iterations}
  $50\!\to\!8$) so $\Delta t$ grows freely when convergence is stable --- verified:
  $\Delta t$ ramps $0.01\!\to\!0.08$ and auto-reduces approaching peak swelling.
\end{itemize}

\section{Roadmap}

\begin{itemize}
  \item \textbf{Phase 2 --- outputs/metrics (ready to start).} Postprocessors for
  contact-area (or delaminated-arc) fraction and the minimum stack pressure to
  maintain contact; multi-cycle loading. No experimental data required.
  \item \textbf{Phase 3 --- design map.} Sweep $\{E\ (\text{or }E/\sigma_y),\,H_v,\,
  G_{Ic}\}$ and produce contour maps of contact loss and required stack
  pressure (the manuscript's opening figure), with a temperature axis via the
  presets.
  \item \textbf{Phase 4 --- DFT integration.} Place DFT-screened candidates on
  the map and recommend the best catholyte.
\end{itemize}

\paragraph{Open calibration items.}
(i)~Recalibrate $E(T),\sigma_y(T),A(T),n$ to the redone hold-segment
nanoindentation at 25/60/90\,$^\circ$C (load--hold--unload--hold protocol;
fused-quartz control; vary hold time and loading rate).
(ii)~Rescale $A$ to the physical loading rate.
(iii)~Decide whether to keep the swelling eigenstrain proxy or move to an
explicit cyclic CAM-state driver.

\appendix
\section{File inventory}
\begin{table}[h]
\centering
\begin{tabular}{ll}
\toprule
File & Role \\
\midrule
\code{calculation/make_mesh.py}   & gmsh generator (conformal / \code{--contact}) \\
\code{calculation/1_contact.i}    & penalty-contact model (J2 catholyte, elastic NVP) \\
\code{calculation/2_czm.i}        & cohesive-zone model (J2 catholyte, elastic NVP) \\
\code{calculation/viscoplastic/}  & parked viscoplastic-catholyte variants (4 inputs) \\
\code{calculation/1_sweep.sh}     & SLURM sweep: $E\times H_v$ \\
\code{calculation/2_sweep.sh}     & SLURM sweep: $E\times H_v\times$ cohesion \\
\code{reference/}                 & NACS manuscript + SI; Shi \textit{et al.} (Matter 2026) \\
\bottomrule
\end{tabular}
\end{table}

\noindent Executable (from \code{calculation/}):
\code{../../../../ecm_test/ecm-opt}.

\end{document}
